% !PRINTED 357
\weberpage{357}

% [sic] source prints $30^0\,40'$; the surrounding calculation implies $33^0\,40'$
Hierauf berechnet man die Differenz (8) mit etwas grösserer Genauigkeit, etwa auf fünf Stellen für einige Winkel in der Nähe der Werthe $30^0\,40'$, von Minute zu Minute fortschreitend, und findet aus der Vorzeichenänderung, dass $\vartheta$ zwischen
\[
33^0\,41' \text{ und } 33^0\,42'
\]
liegt. Wenn man nun für diese beiden Werthe die Rechnung auf sieben Stellen durchführt, so ergiebt sich wieder durch Interpolation der genauere Werth
\[
\vartheta = 33^0 \quad 41' \quad 20{,}6''.
\]
Aus den Formeln (3) oder (4) sieht man, dass hier $r$ negativ ist, und man erhält $r$ sehr einfach aus einer dieser Gleichungen. Man findet die \gesperrt{Brigg'schen} Logarithmen:
\[
\begin{aligned}
\log(-r) &= 0{,}0266148\\
\log\cos\vartheta &= 9{,}9201547 - 10\\
\log\sin\vartheta &= 9{,}7440458 - 10.
\end{aligned}
\]
Also sind die beiden imaginären Wurzeln:
\[
-0{,}884646 \pm 0{,}589740\,i.
\]
